%%%%%%%%%%%%%%%%%%%%%%%%%%%%%%%%%%%%%%%%%
% Short Sectioned Assignment
% LaTeX Template
% Version 1.0 (5/5/12)
%
% This template has been downloaded from:
% http://www.LaTeXTemplates.com
%
% Original author:
% Frits Wenneker (http://www.howtotex.com)
%
% License:
% CC BY-NC-SA 3.0 (http://creativecommons.org/licenses/by-nc-sa/3.0/)
%
%%%%%%%%%%%%%%%%%%%%%%%%%%%%%%%%%%%%%%%%%

%----------------------------------------------------------------------------------------
%	PACKAGES AND OTHER DOCUMENT CONFIGURATIONS
%----------------------------------------------------------------------------------------

\documentclass[fontsize=11pt]{scrartcl} % A4 paper and 11pt font size
\usepackage{geometry} 
\geometry{ 
	a4paper, 
	total={170mm,257mm}, 
	left=20mm, 
	top=20mm, 
} 
\usepackage[T1]{fontenc} % Use 8-bit encoding that has 256 glyphs
\usepackage{fourier} % Use the Adobe Utopia font for the document - comment this line to return to the LaTeX default
\usepackage[english]{babel} % English language/hyphenation
\usepackage{amsfonts,amsthm,amssymb} % Math packages
\usepackage[fleqn]{mathtools}
\usepackage{caption}
\usepackage{lipsum} % Used for inserting dummy 'Lorem ipsum' text into the template
\usepackage{float}
\restylefloat{table}
\usepackage{color}
\usepackage{sectsty} % Allows customizing section commands
\sectionfont{\centering \normalfont} % Make all sections centered, the default font and small caps
\subsectionfont{\fontsize{14}{15}\centering \normalfont} % Make all sections centered, the default font and small caps
\subsubsectionfont{\underline\itshape\fontsize{13}{15}\raggedright\normalfont}

\usepackage{fancyhdr} % Custom headers and footers
\pagestyle{fancyplain} % Makes all pages in the document conform to the custom headers and footers
\fancyhead{} % No page header - if you want one, create it in the same way as the footers below
\fancyfoot[L]{} % Empty left footer
\fancyfoot[C]{} % Empty center footer
\fancyfoot[R]{\thepage} % Page numbering for right footer
\renewcommand{\headrulewidth}{0pt} % Remove header underlines
\renewcommand{\footrulewidth}{0pt} % Remove footer underlines
\setlength{\headheight}{13.6pt} % Customize the height of the header
\newcommand*{\oldneg}{\mathord{\sim}}

\setlength\parindent{0pt} % Removes all indentation from paragraphs - comment this line for an assignment with lots of text

\newcommand{\combination}[2]{
\begin{pmatrix}#1\\#2\end{pmatrix}
}

%----------------------------------------------------------------------------------------
%	TITLE SECTION
%----------------------------------------------------------------------------------------

\newcommand{\horrule}[1]{\rule{\linewidth}{#1}} % Create horizontal rule command with 1 argument of height

\title{	
\normalfont \normalsize 
\textsc{Discrete Math} \\ [25pt] % Your university, school and/or department name(s)
\horrule{0.5pt} \\[0.2cm] % Thin top horizontal rule
\large Homework 3 \\ % The assignment title
\horrule{2pt} \\[0.2cm] % Thick bottom horizontal rule
}
\author{Michael Introini} % Your name

\date{\normalsize{June 19, 2017}} % Today's date or a custom date

\begin{document}
\maketitle % Print the title
\section*{Chapter 6}
\horrule{0.5pt} \\ % Thin top horizontal rule
\subsection*{Section 6.4}
\subsubsection*{Problem 6}
What is the coefficient of $x^7$ in $(1 + x)^{11}$?\\
The coefficient is $\combination{11}{7}$.
\subsubsection*{Problem 12}
The row of Pascal's triangle containing the binomial coefficients 
$
\begin{pmatrix}
	10\\
	k
\end{pmatrix}
$, $ 0 \le k \le 10$, is:\\
$1\quad10\quad45\quad120\quad210\quad252\quad210\quad120\quad45\quad10\quad1$\\
Use Pascal's identity to produce the row immediately following this row in Pascal's triangle.\\

Solution: $1\quad11\quad55\quad165\quad330\quad462\quad330\quad165\quad55\quad11\quad1$ \\
\\
\resizebox{\textwidth}{!}{
	\begin{tabular}{ccccccccccccccccccccccccc}

	&$\combination{10}{0}$&&$\combination{10}{1}$&&$\combination{10}{2}$&&$\combination{10}{3}$&&$\combination{10}{4}$&&$\combination{10}{5}$&&$\combination{10}{6}$&&$\combination{10}{7}$&&$\combination{10}{8}$&&$\combination{10}{9}$&&$\combination{10}{10}$&\\ 

	$\combination{11}{0}$&&$\combination{11}{1}$&&$\combination{11}{2}$&&$\combination{11}{3}$&&$\combination{11}{4}$&&$\combination{11}{5}$&&$\combination{11}{6}$&&$\combination{11}{7}$&&$\combination{11}{8}$&&$\combination{11}{9}$&&$\combination{11}{10}$&&$\combination{11}{11}$
\end{tabular}
}
\subsubsection*{Problem 28}
\textcolor{red}{
Show that if $n$ is a positive integer, then $\combination{2n}{2}= 2\combination{2}{n} + n^2$\\ 
\textbf{a)} using a combinatorial argument\\
\textbf{b)} using Pascal's identity\\
}
\subsection*{Section 6.5}
\subsubsection*{Problem 10}
A croissant shop has plain croissants, cherry croissants,
chocolate croissants, almond croissants, apple croissants, and broccoli croissants. How many ways are there to choose.\\

There are 6 types of croissants to choose from.

\textbf{b)} three dozen croissants?\\
Using theorem 2:\\ $C(n+r-1,r)$, where $n=36$ and $r=6$, we have $C(6+36-1, 36)$ = $\combination{41}{36} = \frac{41!}{35!5!}$ different ways to choose three dozen croissants.\\
 
\textbf{c)} two dozen croissants with at least two of each kind?\\
Since we have 6 kinds of croissants, and we want 24 with at least  2 of each kind, we're actually left to chose pairs of croissants from our available types. Therefore we have $C(6+12-1, 12) = \combination{17}{12} = \frac{17!}{12!5!}$ ways to choose.\\

\subsubsection*{Problem 14}
How many solutions are there to the equation $x_1 + x_2 + x_3 + x_4 = 17$,
where $x_1,x_2,x_3,$ and $x_4$ are nonnegative integers?\\

$C(17+4-1,4) = \frac{20!}{17!3!}$\\

\subsubsection*{Problem 16}
How many solutions are there to the equation $x_1 + x_2 + x_3 + x_4 + x_5 + x_6 = 29$,
where $x_i, i = 1, 2, 3, 4, 5, 6$, is a nonnegative integer such that://

\textbf{a)} $x_i > 1$ for $i=1,2,3,4,5,6$?
Since $x_i$ is a positive integer, then $x_i - 2$ is also a positive integer in our new set. Therefore, $x_1 - 2 + x_2 - 2 + x_3 - 2 + x_4 - 2 + x_5 -2 + x_6 - 2 = 17$. Then we have the following number of possible solutions to the equation.\\
$C(6+17-1, 17) = \combination{22}{17} = \combination{22}{5} = \frac{22!}{17!5!}$

\textbf{b)} $x_1 \ge 1, x_2 \ge 2, x_3 \ge 3, x_4 \ge 4, x_5 \ge 5,$ and $x_6 \ge 6$\\
A solution to the equation with these constraints corresponds to a choice of 1 item of type 1, 2 of type 2, and so on up to 6 items of type 6, together with a choice of 8  $[29 - (1+2+3+4+5+6)=8]$ additional items of any type. Therefore,

$C(6+8-1,8) = C(13,8) = \frac{13!}{8!5!}$\\

\subsubsection*{Problem 40} 
How many ways are there to travel in $xyzw$ space from the origin $(0, 0, 0, 0)$ to the point $(4, 3, 5, 4)$ by taking steps one unit in the positive $x$, positive $y$, positive $z$, or positive $w$ direction?\\

In order to travel from $(0,0,0,0)$ to each of the points in $xyzw$ space, we have to count the possibilities for each point $xyzw$. Therefore we can say that there are $x^4y^3z^5w^4$ arrangements of points. Since each point is a pair, we have $16$ ways to select each point. In total we get the following number of possible ways:\\
$\frac{16!}{4!3!5!4!}$\\

\section*{Chapter 7}
\subsection*{Section 7.1}
\subsubsection*{Problem 6} 
What is the probability that a card selected at random from a standard deck of 52 cards is an ace or a heart?\\

The probability that a card from the deck is a Heart is:\\ $\frac{13}{52} = \frac{1}{4}$\\

The probability that a card from the deck is an Ace is: \\ $\frac{4}{52} = \frac{1}{13}$\\

Then the probability that a card selected at random is an Ace OR a Heart is\\ $\frac{1}{4} + \frac{1}{13} - \frac{5}{52} \approx 0.23$

\subsubsection*{Problem 8}
What is the probability that a five-card poker hand contains the ace of hearts?
There is only 1 ace of hearts in the deck of 52 cards, this means that there are $\combination{1}{1}$ ways to chose it the ace, and $\combination{52}{4}$ ways to chose the other 4 cards in the hand. Then, the probability of a five-card poker hand containing 1 ace of hearts is:\\
$\frac{\combination{1}{1}\combination{52}{4}}{\combination{52}{5}}$\\

\subsubsection*{Problem 14}
What is the probability that a five-card poker hand contains cards of five different kinds?
The number of ways to chose 5 cards of the 13 kinds is $\combination{13}{5}$. For each of these cards there are 4 suits to choose from, therefore there ar $4^5$ distinct choices. Thus, the probability is:\\
$\frac{\combination{13}{5}4^5}{\combination{52}{5}}$\\
 
 
\subsection*{Section 7.2}
\subsubsection*{Problem 10}
What is the probability of these events when we randomly select a permutation of the 26 lowercase letters of the English alphabet?\\

\textbf{d)} a and b are not next to each other in the permutation.\\
We know that the sum of all probabilities is 1 from 
$\displaystyle\sum_{s\in S} p(s) = 1$\\

Then we can say that the probability of $a$ and $b$ not next to each other is actually the difference between the sum of the probability of all possible outcomes and the probability of $a$ and $b$ next to each other.\\
In other words:\\
$1 - \frac{25!+25!}{26!} = 1 - \frac{50}{26} = 1 - \frac{25}{13} = \frac{12}{13} = 92.3\%$\\


\textbf{e)} $a$ and $z$ are separated by at least 23 letters in the permutation.
In order for $a$ and $z$ to be separated by at least 23 letters, we multiply the minimum number of letters between $a$ and $z$ (3) times the total number of letters possible in between (24) times the letters on either end (2). Therefore, the probability is:\\
$\frac{3!\cdot24!\cdot2!}{26!}$\\

\textbf{f)} z precedes both a and b in the permutation.
In order for z to precede $a$ and $b$, we can think of it as $z$ before every other letter, then the probability is:\\
$\frac{25!}{26!}$\\

\subsubsection*{Problem 12}
Suppose that $E$ and $F$ are events such that $p(E) = 0.8$ and $p(F)=0.6$. Show that $p(E\cup F) \ge 0.8$ and $p(E \cap F ) \ge 0.4$.

Using theorem 2:\\
\begin{align*}
p(E \cup F) &= p(E) + p(F) - p(E \cap F)\\
 			&= 0.8 + 0.6 - p(E \cap F)\\
		 	&= 1.4 - p(E \cap F)\\
 			&\text{  However, the probability of any event occurring cannot be more than 1, thus:}  \\
 			&= 1.4 - p(E \cap F) \le 1\\
            &= \mathbf{p(E \cap F) \ge 0.4}
\end{align*}

Now that we've shown $p(E \cap F) \ge 0.4$, we can plug this value into theorem 2 to show $p(E\cup F) \ge 0.8$.

\begin{align*}
p(E \cup F) &= 1.4 - p(E \cap F) \\
			&= 1.4 - 0.6\\
			&\text{ Largest number $p(E \cap F)$ can be}\\
			&= 0.8\\
\end{align*}
\subsubsection*{Problem 16}
Show that if E and F are independent events, then $\bar E$ and $\bar F$ are also independent events.\\
We know that $E$ and $F$ are independent events if and only if $p(E \cap F) = p(E)p(F)$. We can show that $\bar E$ and $\bar F$ in the following manner:\\
\begin{align*}
p(\bar E \cup \bar F) &= p(\bar E)p(\bar F) &\text{By Demorgan's Law}\\
					  &= (1 - p(E))(1 - p(F)) &\text{Theorem 1 Section 7.1}\\
					  &= 1 - p(E) -p(F) + p(E)p(F)\\
					  &= 1 - (p(E) + p(F) - p(E \cap F))\\
					  &= 1 - p(E \cup F) & \text{Laplace's Definition}\\
\end{align*}

To conclude, we get $p(\bar E \cup \bar F) = 1 - p(E \cup F)$ which are equivalent statements.\\
$p(\bar E \cup \bar F)$ : Everything that is \emph{not} in E \emph{not} in F.\\
$1 - p(E \cup F)$ : The universe $(1)$ - everything that \emph{is} in $E$ and everything that \emph{is} in $F$.

\subsubsection*{Problem 26}
Let E be the event that a randomly generated bit string of length three contains an odd number of 1s, and let F be the event that the string starts with 1. Are E and F independent?\\

For the event $p(E)$ we have: $100, 010, 001, 111$\\
For the event $p(F)$ we have: $100, 101, 110, 111$\\

Because we have 8 bit strings of length 3, it follows that\\
$p(E) = p(F) = 4/8 = 1/2$\\

Additionally, because $(E \cap F) = {100,111}$\\
$p(E \cap F) = 2/8 = 1/4$\\

Since $p(E \cap F) = 1/4 = 1/2 \cdot 1/2 = p(E)p(F)$\\
we can conclude that $E$ and $F$ are independent.

\subsubsection*{Problem 34}
Find each of the following probabilities when $n$ independent Bernoulli trials are carried out with probability of success $p$.\\

\textbf{a)} the probability of no successes\\
The probability of exactly 0 success in $n$ independent Bernoulli trials is:
\begin{align*}
	b(0;n,p) &= c(n,0)p^0q^n\\
	         &= c(n,0)q^n
\end{align*}

\textbf{b)} the probability of at least one success\\
Using part \textbf{a} we can assert that there will be at least one success if it is not the case that there are exactly 0 successes. In other words, when $q^n$ is the probability of 0 successes,\\ $p=1-c(n,0)q^n$ will be the probability of at least one success.\\

\subsection*{Section 7.3}
\subsubsection*{Problem 4}
Suppose that Ann selects a ball by first picking one of two boxes at random and then selecting a ball from this box. The first box contains three orange balls and four black balls, and the second box contains five orange balls and
six black balls. What is the probability that Ann picked a ball from the second box if she has selected an orange ball?\\

Let $p(E) = \text{Probability of Chosing an orange ball}$\\
Let $p(\bar E) = \text{Probability of Chosing a black ball}$\\
Let $p(F) = \text{Probability of Chosing from Box 1}$\\
Let $p(\bar F) = \text{Probability of Chosing from Box 2}$\\

Since the choice is random between two boxes,  $p(F) = p(\bar F) = \frac{1}{2}$\\

Th question states that Ann chose an Orange Ball, thus:\\
The probability of the ball coming from Box 1 is: $p(E|F) = \frac{3}{7}$\\
The probability of the ball coming from Box 2 is: $p(E|\bar F) = \frac{5}{11}$\\

\textbf{By the conditional probability definition:}\\
$p(E|F) = \frac{(E \cap F)}{p(F)}$\\
$p(E \cap F) = p(E|F)p(F) = \frac{3}{7} \cdot \frac{1}{2} = \frac{3}{14}$\\

$p(E|\bar F) = \frac{(E \cap \bar F)}{\bar F}$\\
$p(E \cap \bar F) = p(E|\bar F)p(\bar F) = \frac{5}{11} \cdot \frac{1}{2} =  \frac{5}{22}$\\

\textbf{By Baye's Theorem:}\\
$p(\bar F|E) = \frac{p(E|\bar F)p(\bar F)}{p(E|\bar F)p(\bar F) + p(E|F)p(F)}$\\

$p(\bar F|E) = \frac{\frac{5}{22}}{\frac{5}{22} + \frac{3}{14}} = \frac{\frac{5}{22}}{\frac{69}{77}} = \frac{35}{138} \approx 0.25$\\

\subsubsection*{Problem 6}
\textcolor{red}{When a test for steroids is given to soccer players, 98\% of the players taking steroids test positive and 12\% of the players not taking steroids test positive. Suppose that 5\% of soccer players take steroids. What is the probability that a soccer player who tests positive takes steroids?}\\

\subsection*{Section 7.4}:
\subsubsection*{Problem 6}
What is the expected value when a \$1 lottery ticket is bought in which the purchaser wins exactly \$10 million if the ticket contains the six winning numbers chosen from the set ${1, 2, 3, \cdots 50}$ and the purchaser wins nothing otherwise?\\

We let the probability of choosing the winning combination be $(10^7 - 1)\frac{6!44!}{50!}$.\\ where $10^7 -1$ is the amount of money won after deducting the cost of the ticket.\\
And we let the probability of not choosing the winning combination be $(-1)(1 - \frac{6!44!}{50!})$ where $-1$ is the amount of money lost if we did not win the prize.\\

$E(X) = (10^7 - 1)\frac{6!44!}{50!} + (-1)(1 - \frac{6!44!}{50!}) \approx -0.37 $

The expected value of the ticket is $-\$0.37$\\
\subsection*{Section 8.1}

\subsubsection*{Problem 12}
\textbf{a)} Find a recurrence relation for the number of ways to climb n stairs if the person climbing the stairs can take one, two, or three stairs at a time.\\

The recurrence relation for this sequence is $a_n = a_{n-1} + a_{n-2} + a_{n-3}$ since we can have up to $n-3$ ways of climbing steps of size $n$.\\

\textbf{b)} What are the initial conditions?\\
Thus, the Initial conditions are:\\
$a_0 = 1, a_1 = 1, a_2 = 2, a_3 = 4$\\

\textbf{Because there is:}\\
1 way to climb 0 steps\\
1 way to climb 1 step\\
2 ways to climb 2 steps\\
- 1 step, then 1 step\\
- 2 steps at once\\
4 ways to climb 3 steps\\
- 1 step, 1 step, 1 step\\
- 1 step, 2 steps at once\\
- 2 steps at once, 1 step\\
- 3 steps at once\\ 
\textbf{c)} In how many ways can this person climb a flight of eight stairs?\\ 

First we need to find $a_4, a_5, a_6, a_7$\\
$a_4 = 7$, $a_5 = 13$, $a_6 = 24$, $a_7 = 44$\\

The total is:\\
$a_8 = a_{8-1} + a_{8-2} + a_{8-3} = a_{7} + a_{6} + a_{5} = 81$\\

\subsubsection*{Problem 26}
\textbf{a)} Find a recurrence relation for the number of ways to completely cover a 2 × n checkerboard with 1 x 2 dominoes. [Hint: Consider separately the coverings where the position in the top right corner of the checkerboard is covered by a domino positioned horizontally and where it is covered by a domino positioned vertically.]\\

The recurrence relation is $a_n = a_{n-1} + a_{n-2}$

\textbf{b)} What are the initial conditions for the recurrence relation in part (a)?\\

The initial condition is $a_0 = 1$, $a_1 = 1$, $a_2 = 3$, $a_3 = 4$\\

\textbf{c)} How many ways are there to completely cover a 2 × 17 checkerboard with 1 × 2 dominoes?\\

The sum is $1597$.\\

\subsection*{Section 8.2}
\subsubsection*{Problem 4}
Solve these recurrence relations together with the initial conditions given.\\

\textbf{b)} $a_n = 7a_{n-1} - 10a_{n-2}$ for $n \ge 2, a_0 = 2 , a_1 = 1$\\

Let $x^n$ a variable\\
$x^n = 7a_{n-1} - 10a_{n-2}$\\
$x^n -7a^{n-1} + 10a^{n-2} = 0$\\
$x^n (-7a^{n-1} + 10a^{n-2}) = 0$\\
$x^2 (-7a + 10) = 0$\\
$x^2  (a - 5)(a - 2) = 0$\\

Then the sequence $a_n$ is a solution to the recurrence relation if and only if there exists and equation such that,\\
$a_1 2^n + a_2 5^n$ where $a_1, a_2$ are constants.\\

From the initial condition, we have:\\
$a_0 = 2 = a_1 + a_2$\\
$a_1 = 1 = 2a_1 + 5a_2$

Solving each we get:\\
$a_1 = 3$ and $a_2 = -1$\\

Which leaves us with the solution $a_n = 3\cdot2^n -1\cdot5^n$

\end{document} 
